
\section{Broader Related Work}
\label{app:related}

Compared to a concurrent work, Reflexion \cite{shinn2023reflexion}, our approach involves correction using feedback, whereas their setup involves finding the next best solution in planning using ReAct. While ReAct and Reflexion provide a free-form reflection on whether a step was executed correctly and potential improvements, our approach is more granular and structured, with multi-dimensional feedback and scores. This distinction allows our method to offer more precise and actionable feedback, making it suitable for a wider range of natural language generation tasks, including those that may not necessarily involve step-by-step planning such as open-ended dialogue generation. 


\paragraph{Comparison with \citet{Welleck2022SelfCorrect}} 
The closest work to ours may be Self-Correction \citep{Welleck2022SelfCorrect}; however, Self-Correction has several disadvantages compared to \ours:
\begin{enumerate}[itemsep=0pt,topsep=0pt,partopsep=0pt]
    \item Self-Correction does not train their model to generate explicit feedback; instead, \citet{Welleck2022SelfCorrect} trained their models to refine only. As we show in \Cref{sec:analysis} and \Cref{tab:ablation_feedback}, having the model generate explicit feedback results in significantly better refined outputs.
    \item Self-Correction trains a separate refiner (or ``corrector'') for each task. In contrast, \ours uses instructions and few-shot prompting, and thus does not require training a separate refiner for each task.
    \item Empirically, we evaluated \ours using the same base model of GPT-3 as Self-Correction, and with the same settings on the GSM8K benchmark. Self-Correction  
    achieved 45.9\% accuracy while \ours (this work) achieved \textbf{55.7}\% ($\bm{\uparrow}$\textbf{9.8}).
\end{enumerate}


\paragraph{Comparison with non-refinement reinforcement learning (RL) approaches.}
Rather than having an explicit refinement module, an alternative way to incorporate feedback is by optimizing a scalar reward function, e.g. with reinforcement learning (e.g.,~\citet{stiennon2020learning,Lu2022QuarkCT,Le2022CodeRL}).
These methods differ from \ours (and more generally, refinement-based approaches) in that the model cannot access feedback on an intermediate generation. 
Second, these reinforcement learning methods require updating the model's parameters, unlike \ours.




See Table \ref{tab:relatedwork} for an additional detailed comparison of related work.


\begin{table}[H]
\centering
\begin{adjustbox}{max width=\textwidth}
\scriptsize
\begin{tabular}{lp{3cm}llp{2cm}l}
\toprule
 \textbf{Method}&
  \textbf{Primary Novelty} &
  zero/few shot improvement &
  multi aspect critics &
  NL feedback with error localization &
  iterative framework \\ \midrule
\cellcolor[HTML]{CBCEFB}RLHF \cite{stiennon2020learning} &
  optimize for human preference &
  \emoji[emoji]{x}trained on feedback &
  \emoji[emoji]{x} single (human) &
  \emoji[emoji]{check}(not self gen.) &
  \emoji[emoji]{x} \\
\cellcolor[HTML]{CBCEFB}Rainier RL \cite{Liu2022RainierRK} &
  RL to generate knowledge &
  \emoji[emoji]{x}trained on end task &
  \emoji[emoji]{x} single(accuracy) &
  \emoji[emoji]{x}(knowl. only) &
  \emoji[emoji]{x} \\
\cellcolor[HTML]{CBCEFB}\textsc{Quark} RL \cite{Lu2022QuarkCT} &
  quantization to edit generations &
  \emoji[emoji]{x} trained on end task &
  \emoji[emoji]{x} single(scalar score) &
  \emoji[emoji]{x} (dense signal) &
  \emoji[emoji]{check} (train time iter.) \\
\cellcolor[HTML]{CBCEFB}Code RL \cite{Le2022CodeRL} &
  actor critic RL for code improvement &
  \emoji[emoji]{x} trained on end task &
  \emoji[emoji]{x} single(unit tests) &
  \emoji[emoji]{x}(dense signal) &
  \emoji[emoji]{x} \\ \midrule
\cellcolor[HTML]{FFCC67}DrRepair \cite{Yasunaga2020DrRepair} &
  Compiler feedback to iteratively repair &
  \emoji[emoji]{x} trained semi sup. &
  \emoji[emoji]{x} single(compiler msg) &
  \emoji[emoji]{check}(not self gen.) &
  \emoji[emoji]{check} \\
\cellcolor[HTML]{FFCC67}PEER \cite{Schick2022PEERAC} &
  doc. edit trained on wiki edits &
  \emoji[emoji]{x}trained on edits &
  \emoji[emoji]{x} single(accuracy) &
  \emoji[emoji]{check}(not self gen.) &
  \emoji[emoji]{check} \\
\cellcolor[HTML]{FFCC67}Self critique \cite{Saunders2022SelfCritique} &
  few shot critique generation &
  \emoji[emoji]{x}feedback training &
  \emoji[emoji]{x} single(human) &
  \emoji[emoji]{check}(self gen.) &
  \emoji[emoji]{x} \\ 
\cellcolor[HTML]{FFCC67}Self-correct \cite{Welleck2022SelfCorrect} &
  novel training of a corrector &
  \emoji[emoji]{x} trained on end task &
  \emoji[emoji]{x} single (task specific) &
  \emoji[emoji]{check}(limited setting) &
  \emoji[emoji]{check}(limited setting) \\ 
  \cellcolor[HTML]{FFCC67}Const. AI \cite{bai2022constitutional} &
  train RL4F on automat (critique, revision) pair &
  \emoji[emoji]{x} critique training &
  \emoji[emoji]{check} (fixed set) &
  \emoji[emoji]{check} &
  \emoji[emoji]{x} \\ \midrule
\cellcolor[HTML]{9AFF99}Self-ask \cite{press2022measuring} &
  ask followup ques when interim ans correct;final wrong &
  \emoji[emoji]{check} few shot &
  \emoji[emoji]{x} none &
  \emoji[emoji]{x}(none) &
  \emoji[emoji]{x} \\
\cellcolor[HTML]{9AFF99}GPT3 score \cite{fu2023gptscore} &
  GPT can score generations with instruction &
  \emoji[emoji]{check}few shot &
  \emoji[emoji]{x} single(single utility fn) &
  \emoji[emoji]{x}(none) &
  \emoji[emoji]{x} \\
\cellcolor[HTML]{9AFF99}Augmenter \cite{Peng2023LLMAugmenter} &
  factuality feedback from external KBs &
  \emoji[emoji]{check}few shot &
  \emoji[emoji]{x} single(factuality) &
  \emoji[emoji]{check}(self gen.) &
  \emoji[emoji]{check} \\
\cellcolor[HTML]{9AFF99}Re$^3$ \cite{Yang2022Re3GL} &
  $\sim$ours: but one domain, trained critics &
  \emoji[emoji]{check}few shot &
  \emoji[emoji]{check}(trained critics) &
  \emoji[emoji]{check}(not self gen.) &
  \emoji[emoji]{check} \\
\cellcolor[HTML]{9AFF99}\ours &
  fewshot iterative multi aspect NL fb &
  \emoji[emoji]{check}few shot &
  \emoji[emoji]{check} multiple(few shot critics) &
  \emoji[emoji]{check}(self gen.) &
  \emoji[emoji]{check}
\end{tabular}
\end{adjustbox}
\normalsize
\caption{Summary of related approaches. Reinforcement learning approaches are shown in \colorbox[HTML]{CBCEFB}{purple}}, trained corrector approaches are shown in \colorbox[HTML]{FFCC67}{orange}, and few-shot corrector approaches are shown in \colorbox[HTML]{9AFF99}{green}.
\label{tab:relatedwork}
\end{table}

